\documentclass[11pt]{article}
\usepackage[a4paper,margin=27mm]{geometry}
\usepackage{amsmath,amssymb,amsthm}
\usepackage{graphicx}
\usepackage{microtype}
\usepackage[hidelinks]{hyperref}

\newtheorem{theorem}{Theorem}
\newtheorem{lemma}{Lemma}
\newcommand{\R}{\mathbb R}
\newcommand{\C}{\mathbb C}
\newcommand{\Sone}{\boldsymbol S_1}
\newcommand{\Tr}{\operatorname{Tr}}
\newcommand{\diver}{\operatorname{div}}

\title{An absolutely continuous trace formula for\\strongly commuting normal operators}
\author{Run \texttt{fa\_banach\_001}, lane 9\\\texttt{agent\_lane\_09}}
\date{11 August 2026}

\begin{document}
\maketitle

\begin{center}
\fbox{\parbox{0.91\textwidth}{\textbf{Status.}
Candidate partial result, likely valid.  It answers both the measure-existence
and absolute-continuity parts of the source problem under the additional
hypothesis that the two normal operators commute with one another.}}
\end{center}

\section{Source question}

Aleksandrov, Peller, and Potapov~\cite{APP} ask the following in Open
Problem~2 on page~6.  If normal operators \(N_1,N_2\) have trace-class
difference, do finite Borel measures \(\nu_1,\nu_2\) on \(\R^2\) exist such
that
\begin{equation}\label{eq:source}
 \Tr\bigl(f(N_2)-f(N_1)\bigr)
 =\int_{\R^2}f_x\,d\nu_1+\int_{\R^2}f_y\,d\nu_2
\end{equation}
for sufficiently smooth \(f\)?  If so, may the measures be chosen absolutely
continuous with respect to planar Lebesgue measure?

\clearpage
\begin{figure}[!ht]
\centering
\includegraphics[width=0.98\textwidth,height=0.78\textheight,keepaspectratio]{figures/open_problem_crop.png}
\caption{Open Problem 2 in arXiv:1901.09495, page 6.}
\label{fig:question}
\end{figure}
\clearpage

The theorem below gives an affirmative answer to both questions when
\(N_1N_2=N_2N_1\).  This is stronger than merely asking that the real and
imaginary parts of each individual \(N_j\) commute; that latter condition is
automatic from normality and is part of the unrestricted problem.

\section{Proof intuition}

Cross-commutation makes the trace-class perturbation
\(K=N_2-N_1\) a compact normal operator commuting with \(N_1\).  Every
nonzero eigenspace of \(K\) is finite-dimensional and reduces \(N_1\).  On the
\(\kappa\)-eigenspace, diagonalising \(N_1\) turns \(N_2\) into the same
diagonal matrix shifted by \(\kappa\).  The entire trace is therefore a
summable collection of scalar endpoint differences
\(f(b)-f(a)\), whose total travel length is exactly \(\|K\|_{\Sone}\).

The obvious representation integrates \(\nabla f\) along the segment from
\(a\) to \(b\), but its measure is singular.  To get absolute continuity,
replace the segment by a compactly supported planar flow.  First spread each
endpoint to a smooth blob at scale \(|b-a|\), using a radial field with one
point source.  Then transport one blob to the other along the segment.  The
resulting \(L^1\) vector field \(V_{a,b}\) has
\(\diver V_{a,b}=\delta_a-\delta_b\) and
\(\|V_{a,b}\|_1\leq3|b-a|\).  Summability in trace norm now permits all these
absolutely continuous flows to be added.

\section{A planar endpoint field}

We identify \(\C\) with \(\R^2\), write \(m_2\) for planar Lebesgue measure,
and use the Euclidean norm for vector fields.

\begin{lemma}\label{lem:field}
For every \(a,b\in\R^2\), there is a compactly supported real vector field
\(V_{a,b}\in L^1(\R^2;\R^2)\) such that
\[
 \diver V_{a,b}=\delta_a-\delta_b
 \quad\text{in distributions},
 \qquad
 \|V_{a,b}\|_{L^1}\leq3|b-a|.
\]
Consequently, for every \(f\in C^1(\R^2)\),
\begin{equation}\label{eq:endpoint}
 f(b)-f(a)=\int_{\R^2}\nabla f(x)\mathbin{\cdot}V_{a,b}(x)\,dm_2(x).
\end{equation}
\end{lemma}

\begin{proof}
The case \(a=b\) is immediate.  Fix a nonnegative radial
\(\eta\in C_c^\infty(B(0,1))\) with \(\int\eta\,dm_2=1\), and set
\[
 M(r)=\int_{|x|\leq r}\eta(x)\,dm_2(x),
 \qquad
 G(x)=\frac{1-M(|x|)}{2\pi}\frac{x}{|x|^2}\quad(x\ne0).
\]
Since \(M(r)=1\) for \(r\geq1\), the field \(G\) is compactly supported.
Radial integration gives
\[
 \|G\|_1=\int_0^1(1-M(r))\,dr\leq1.
\]
The flux of \(G\) through a small circle about zero tends to one, while away
from zero the radial divergence formula gives \(\diver G=-\eta\).  Hence
\begin{equation}\label{eq:Gdiv}
 \diver G=\delta_0-\eta
\end{equation}
distributionally.

Put \(h=b-a\), \(r=|h|\),
\(\eta_r(x)=r^{-2}\eta(x/r)\), and \(G_r(x)=r^{-1}G(x/r)\).  Define
\begin{align*}
 D(x)&=G_r(x-a)-G_r(x-b),\\
 H(x)&=h\int_0^1\eta_r(x-a-th)\,dt,\\
 V_{a,b}(x)&=D(x)+H(x).
\end{align*}
Scaling \eqref{eq:Gdiv} yields
\[
 \diver D=\delta_a-\eta_r(\,\cdot-a)-\delta_b+\eta_r(\,\cdot-b).
\]
Moreover,
\[
 \diver H
 =\int_0^1 h\mathbin{\cdot}\nabla\eta_r(x-a-th)\,dt
 =\eta_r(x-a)-\eta_r(x-b).
\]
Thus \(\diver V_{a,b}=\delta_a-\delta_b\).  Both fields are compactly
supported, and
\[
 \|D\|_1\leq2r\|G\|_1\leq2r,
 \qquad
 \|H\|_1\leq |h|\int_0^1\|\eta_r\|_1\,dt=r.
\]
This proves the norm estimate.  Finally, distributional integration by parts
gives
\(-\int V_{a,b}\cdot\nabla f=f(a)-f(b)\), which is
\eqref{eq:endpoint}.  A cutoff equal to one on the compact support of the
field justifies the same identity when \(f\) is not compactly supported.
\end{proof}

\section{The partial trace formula}

\begin{theorem}\label{thm:main}
Let \(N_1,N_2\) be bounded normal operators on a complex Hilbert space.
Assume
\[
 N_1N_2=N_2N_1,
 \qquad K:=N_2-N_1\in\Sone.
\]
Then there is a real vector field \(v=(v_1,v_2)\in
L^1(\R^2;\R^2)\) such that
\begin{equation}\label{eq:main}
 \Tr\bigl(f(N_2)-f(N_1)\bigr)
 =\int_{\R^2}\bigl(f_xv_1+f_yv_2\bigr)\,dm_2
\end{equation}
for every complex-valued \(f\in C^1(\R^2)\).  The operator difference in
\eqref{eq:main} is trace class, and the densities satisfy
\begin{equation}\label{eq:norms}
 \int_{\R^2}|v|\,dm_2\leq3\|K\|_{\Sone},
 \qquad
 \|v_1\|_1+\|v_2\|_1
 \leq3\sqrt2\,\|K\|_{\Sone}.
\end{equation}
In particular, \(\nu_j=v_jm_2\) are finite real signed Borel measures,
absolutely continuous with respect to \(m_2\), and they solve the source trace
formula in this subcase.
\end{theorem}

\begin{proof}
Because commuting bounded normal operators doubly commute by Fuglede's
theorem~\cite{Fuglede}, \(K\) is normal and commutes with both \(N_1\) and
\(N_1^*\).  Let \(\{\kappa\}\) be the nonzero eigenvalues of the compact
normal operator \(K\), and let \(E_\kappa=\ker(K-\kappa I)\).  These spaces
are finite-dimensional, mutually orthogonal, and reduce \(N_1\).  Their
orthogonal complement is \(\ker K\), where \(N_2=N_1\).

On each \(E_\kappa\), choose an orthonormal basis diagonalising the normal
restriction of \(N_1\); write its eigenvalues as
\(\lambda_{\kappa,1},\ldots,\lambda_{\kappa,m_\kappa}\).  In the same basis,
the eigenvalues of \(N_2\) are
\(\lambda_{\kappa,j}+\kappa\).  Trace-class summability says
\begin{equation}\label{eq:sumk}
 \sum_{\kappa\ne0}m_\kappa|\kappa|=\|K\|_{\Sone}.
\end{equation}

For every pair
\(a=\lambda_{\kappa,j}\), \(b=\lambda_{\kappa,j}+\kappa\), take the field
from Lemma~\ref{lem:field}, and set
\[
 v=\sum_{\kappa\ne0}\sum_{j=1}^{m_\kappa}
 V_{\lambda_{\kappa,j},\lambda_{\kappa,j}+\kappa}.
\]
By the lemma and \eqref{eq:sumk}, this series converges absolutely in
\(L^1(\R^2;\R^2)\), and its norm is at most \(3\|K\|_{\Sone}\).  The second
bound in \eqref{eq:norms} follows from
\(|v_1|+|v_2|\leq\sqrt2|v|\).

All endpoints and field supports lie in a fixed compact set depending only on
\(\|N_1\|\) and \(\|K\|\).  Hence \(\nabla f\) is bounded there.  The scalar
mean-value estimate and \eqref{eq:sumk} show that the block diagonal operator
\(f(N_2)-f(N_1)\) is trace class.  Absolute convergence, followed by
Lemma~\ref{lem:field}, now gives
\begin{align*}
 \Tr\bigl(f(N_2)-f(N_1)\bigr)
 &=\sum_{\kappa\ne0}\sum_{j=1}^{m_\kappa}
 \bigl[f(\lambda_{\kappa,j}+\kappa)-f(\lambda_{\kappa,j})\bigr]\\
 &=\int_{\R^2}\nabla f\mathbin{\cdot}v\,dm_2,
\end{align*}
which is \eqref{eq:main}.
\end{proof}

\section{Verification, limitations, and novelty status}

The proof is deterministic and has no computational dependency.  The two
points most worth checking are the sign and scaling in
\(\diver V_{a,b}=\delta_a-\delta_b\), and the use of Fuglede's theorem to make
each nonzero eigenspace of \(K\) reducing for \(N_1\).  The construction also
works on nonseparable Hilbert spaces: the trace-class range is separable and
the complement lies in \(\ker K\).

The cross-commutation hypothesis is a serious limitation.  For arbitrary
normal matrices \(A,B\), the best possible trace bound over real
one-Lipschitz \(f\) is the unnormalised planar \(W_1\) matching cost between
their eigenvalue multisets.  Thus a dimension-free estimate of that cost by
\(\|A-B\|_{\Sone}\) would provide the missing finite-dimensional input; an
unbounded sequence of ratios could instead feed a direct-sum counterexample.
Hoffman--Wielandt controls \(W_2\) by the Hilbert--Schmidt norm, but its
conversion to \(W_1\) loses a factor \(\sqrt n\).  The present proof does not
resolve this hinge.

A bounded literature search through 11 August 2026 covered the exact problem
wording, normal-operator trace formulas, commuting self-adjoint tuples,
spectral-shift measures in two variables, and normal-matrix matching in
Schatten norms.  Chattopadhyay--Sinha~\cite{CS} prove a two-variable
Stokes-like formula for a second-order expression under Hilbert--Schmidt
perturbations; it is not the first-order formula \eqref{eq:source}.
Aleksandrov--Peller--Potapov--Sukochev~\cite{APPS} prove Besov/Schatten
estimates for normal perturbations, but not the requested Lipschitz trace
estimate.  No exact statement of Theorem~\ref{thm:main} was found.  Novelty
confidence is moderate rather than high because the argument becomes
elementary after imposing cross-commutation.

\paragraph{Human-review recommendation.}
Retain this as a substantial partial theorem if the explicit \(L^1\) endpoint
field and the block trace argument verify.  Do not treat it as resolving the
arbitrary-normal question.  The most valuable next step is the
\(\Sone\)-to-\(W_1\) spectral matching hinge identified above.

\begin{thebibliography}{9}
\bibitem{APP}
A. B. Aleksandrov, V. V. Peller, and D. S. Potapov,
\emph{On a trace formula for functions of noncommuting operators},
arXiv:1901.09495 (2019).

\bibitem{APPS}
A. B. Aleksandrov, V. V. Peller, D. Potapov, and F. Sukochev,
\emph{Functions of normal operators under perturbations},
Adv. Math. 226 (2011), 5216--5251; arXiv:1008.1638.

\bibitem{CS}
A. Chattopadhyay and K. B. Sinha,
\emph{Trace formula for two variables}, arXiv:1402.0792 (2014).

\bibitem{Fuglede}
B. Fuglede,
\emph{A commutativity theorem for normal operators},
Proc. Natl. Acad. Sci. USA 36 (1950), 35--40.
\end{thebibliography}

\end{document}
